% SEsml: Schema-guided Extraction with Small Language Models for GenSIE 2026
% CEURART 1-column format (ceurart.cls v0.62+)
% Proceedings of IberLEF 2026 @ SEPLN

\documentclass[onecolumn]{ceurart}

\usepackage[utf8]{inputenc}
\usepackage[T1]{fontenc}
\usepackage{libertinus}
\usepackage{booktabs}
\usepackage{graphicx}
\usepackage{xurl}
\usepackage{hyperref}

% Bibliography
\usepackage[authoryear,round]{natbib}
\bibliographystyle{plainnat}

% Metadata
\def\Year{2026}
\def\Publisher{CEUR-WS.org}

\begin{document}

% ==============================================================================
% FRONT MATTER
% ==============================================================================

\title{SEsml: Schema-guided Extraction with Small Language Models for GenSIE 2026}

\author[1]{Alejandro Beltr\'{a}n}[%
orcid=0009-0007-9069-2039,
email=alejandro.beltran@rect.uh.cu,
]
\author[1]{Daniel Toledo}[%
email=daniel.toledo@matcom.uh.cu,
]

\address[1]{Universidad de La Habana, Havana, Cuba}

\abstract{
  We present SEsml, a system for zero-shot schema-guided information extraction
  from Spanish text using Small Language Models (SLMs) under 14 billion
  parameters, developed for the GenSIE 2026 shared task at IberLEF 2026. Our
  core contribution is a \textbf{draft-then-decode} architecture that explicitly
  decouples semantic content extraction from structural JSON formatting across
  two sequential LLM calls. The first phase generates a semi-structured
  \texttt{campo: valor} draft that captures extracted values without syntactic
  overhead, mitigating the \emph{cognitive collapse} that SLMs experience when
  simultaneously reasoning about Spanish linguistic content and enforcing JSON
  grammar. The second phase maps the draft to a strictly valid JSON object
  through constrained decoding (\texttt{response\_format=json\_schema}), with a
  local fallback parser for robustness. We submitted three pipelines: the
  principal two-phase extraction pipeline, a lightweight single-call
  HybridCoT pipeline with visible Chain-of-Thought reasoning, and an adaptive
  pipeline that dynamically assembles prompts via semantic similarity to
  precomputed development-set profiles. Our SchemaOptimizer preprocessing
  module normalizes arbitrary JSON Schemas --- inlining \$ref references,
  stripping unsafe keys, and enforcing \texttt{additionalProperties: false} ---
  reducing constrained-decoding engine failures by an estimated 15\% and
  grammar compilation time by 3--5$\times$. At the time of writing, official
  evaluation results are pending; we present the architectural design, prompt
  engineering methodology, error analysis, and a comprehensive plan for the
  forthcoming experimental evaluation.
}

\begin{keywords}
  Schema-guided information extraction \sep
  Small language models \sep
  Constrained decoding \sep
  Chain-of-thought \sep
  Spanish NLP \sep
  Zero-shot extraction
\end{keywords}

\maketitle

% ==============================================================================
% 1. INTRODUCTION
% ==============================================================================

\section{Introduction}
\label{sec:introduction}

GenSIE 2026 (General-purpose Schema-guided Information Extraction), held as
part of IberLEF 2026, introduced a challenging shared task: given a fragment of
Spanish text, a natural-language instruction, and an arbitrary JSON Schema
provided at inference time, systems must produce a valid JSON object whose
structure adheres exactly to the schema and whose values are faithfully
extracted from the source text. The task is zero-shot --- no per-schema training
or few-shot examples are permitted --- and targets Small Language Models (SLMs)
with fewer than 14 billion parameters, reflecting a growing interest in
efficient, sovereign NLP deployments.

Extracting structured JSON under these constraints is fundamentally harder than
open-ended text generation. SLMs, when asked to produce JSON directly,
frequently exhibit what we term \emph{cognitive collapse}: the model's limited
representational capacity is torn between reasoning about linguistic
content --- particularly in a morphologically rich language like Spanish --- and
simultaneously enforcing the rigid syntactic conventions of JSON (brackets,
commas, quotation marks, strict typing). This dual burden leads to a
characteristic pattern of failures: truncated output, misplaced brackets,
omitted fields, and hallucinated values that fill structural gaps demanded by
the schema.

We address this challenge with a \textbf{two-phase draft-then-decode}
architecture that explicitly decouples content reasoning from formatting. This
approach is related to the Generate \& Organize (G\&O) paradigm proposed
by~\citet{li2024generate} for structured text generation, and is theoretically
grounded in the analysis of~\citet{reddy2026dccd}, who formalise the
\emph{projection tax} incurred by standard constrained decoding when the
language model assigns low probability mass to valid continuations. Rather than
asking the SLM to produce valid JSON in a single step, our principal pipeline
first generates a semi-structured draft in a simplified \texttt{campo: valor}
notation, where the model focuses solely on identifying and extracting the
relevant spans of information. A second, constrained decoding step then maps
these extracted values into the exact target JSON schema using
\texttt{response\_format=json\_schema}, making syntactic violations physically
impossible. A local fallback parser ensures robustness even when the structured
output API fails.

We submitted three pipelines to the GenSIE 2026 evaluation. The
\textbf{extraction} pipeline, our principal submission, implements the full
two-phase approach with two LLM calls. The \textbf{hybrid\_cot} pipeline is a
lightweight single-call variant that interleaves visible Chain-of-Thought
reasoning~\citep{wei2022chain} with direct JSON generation using an internal
\texttt{\_extract\_json} heuristic. The \textbf{adaptive} pipeline employs a
single LLM call whose prompt is assembled dynamically from modular components
selected via semantic similarity between the current task and precomputed
development-set vectors.

Our team, \textbf{SEsml (Schema-guided Extraction with Small Language Models)},
consists of Alejandro Beltr\'{a}n and Daniel Toledo from the Universidad de La
Habana. All pipelines are implemented in Python, containerised via Docker, and
exposed as a FastAPI server compatible with the GenSIE evaluation harness. The
code is publicly available\footnote{\url{https://github.com/alejbv/SEsml.git}}.

% ==============================================================================
% 2. SYSTEM ARCHITECTURE
% ==============================================================================

\section{System Architecture}
\label{sec:architecture}

All three pipelines share a common architectural philosophy: minimise the
cognitive burden on the SLM by delegating structural enforcement to
deterministic components wherever possible. Every pipeline begins with the
\textbf{SchemaOptimizer} (described in detail in Section~\ref{sec:schema-opt}),
a local, LLM-free module that preprocesses the incoming JSON Schema into a form
that is safe and efficient for constrained decoding engines. From this shared
starting point, the pipelines diverge in their allocation of LLM calls and
their strategy for bridging extracted content to structured output.

Table~\ref{tab:pipelines} summarises the three pipelines at a glance.

\begin{table}[h]
\centering
\caption{Overview of the three SEsml pipelines submitted to GenSIE 2026.}
\label{tab:pipelines}
\begin{tabular}{lll}
\toprule
\textbf{Pipeline} & \textbf{LLM Calls} & \textbf{Approach} \\
\midrule
extraction & 2 & SchemaOptimizer $\rightarrow$ DraftEngine \\
           &   & (\texttt{campo: valor}) $\rightarrow$ Constrained Decoding \\
           &   & $\rightarrow$ Local fallback parser \\
\addlinespace
hybrid\_cot & 1 & SchemaOptimizer $\rightarrow$ Single call \\
            &   & with visible CoT $\rightarrow$ Direct JSON \\
\addlinespace
adaptive    & 1 & SchemaAnalyzer $\rightarrow$ PromptAssembler \\
            &   & $\rightarrow$ Single call with adaptive prompt \\
\bottomrule
\end{tabular}
\end{table}

\subsection{Principal Pipeline: Extraction}

The extraction pipeline operates in two distinct LLM calls. The first call,
managed by the \textbf{DraftEngine}, receives the optimised schema and the
source text and is instructed to produce a draft in \texttt{campo: valor}
format --- a simple line-oriented notation where each extracted field is written
as its path followed by a colon and its value. The draft is wrapped in
\texttt{<Thinking>...</Thinking>} and \texttt{<Draft>...</Draft>} XML tags that
provide explicit structural boundaries. This format eliminates all JSON
syntactic overhead: the model writes no brackets, no commas, and no quotation
marks around keys, substantially reducing the risk of cognitive collapse.

The prompt template for this stage (version~v3b, approximately 1,163 characters
or 387 tokens) encodes three critical rule sets. The \texttt{[INFERENCIA]} rule
states that the model may only perform semantic inference when the schema's
field \texttt{description} contains one of the trigger words \texttt{reasoned},
\texttt{inferred}, \texttt{based on}, or \texttt{classify}. The
\texttt{[VALORES]} rule demands exact string matching, case-sensitive enum
adherence, and mandates \texttt{null} when information is absent. The
\texttt{[FORMATO]} rule instructs the model to follow each field's
\texttt{description} for formatting guidance rather than applying fixed global
rules.

The second LLM call, the \textbf{Decoder}, takes the optimised schema and the
draft from the first stage and maps the extracted values into a strictly valid
JSON object using \texttt{response\_format=json\_schema}. Because the decoder
operates on an already-extracted draft, its task is purely one of formatting and
does not require re-reasoning about the source text. If the structured output
API returns a malformed response, a local fallback parser extracts values from
the draft directly using regular-expression based heuristics over the
\texttt{campo: valor} lines.

\subsection{Alternative Pipelines}

The \textbf{hybrid\_cot} pipeline collapses the two-phase process into a single
LLM call. The model is prompted to reason aloud in a \texttt{<Thinking>}
block --- producing a visible Chain-of-Thought~\citep{wei2022chain} trace ---
and then to output the final JSON directly. A lightweight
\texttt{\_extract\_json} function extracts the JSON object from the model's
response using heuristics that locate the first \texttt{\{} and the last
\texttt{\}} brace. This pipeline trades the accuracy gains of constrained
decoding for lower latency and reduced token consumption.

The \textbf{adaptive} pipeline introduces a more sophisticated prompt
construction mechanism. Before the single LLM call, a \textbf{SchemaAnalyzer}
computes a task profile consisting of schema complexity metrics, field-type
distributions, domain cues, and known error patterns. A
\textbf{PromptAssembler} then selects and orders modular prompt components
within a strict 1,400-character budget, guided by cosine similarity between the
embedding of the current task's instruction (computed via fastembed using
BAAI/bge-small-en-v1.5) and precomputed embeddings of the development-set
tasks, following a retrieval-augmented prompt construction strategy.

\subsection{Common Infrastructure}

All three pipelines share a uniform runtime infrastructure. Inference is
performed through an OpenAI-compatible client that accepts any backend exposing
the chat completions interface. The \texttt{model} parameter is passed through
transparently from the GenSIE evaluator, making the system model-agnostic. All
calls use \texttt{temperature=0.0} and disable streaming to ensure
deterministic, complete responses. Token usage is tracked by a shared
\texttt{UsageTracker} module.

The system is packaged as a Docker container that serves a FastAPI application
on port~8000. Two endpoints are exposed: \texttt{/info}, which returns the
participant identifier and a list of registered pipeline names, and
\texttt{/run}, which accepts the task payload and returns the extracted JSON
together with metadata. This design conforms to the GenSIE evaluation harness
specification.

% ==============================================================================
% 3. SCHEMA OPTIMIZATION
% ==============================================================================

\section{Schema Optimization}
\label{sec:schema-opt}

The \textbf{SchemaOptimizer} is the entry point for every pipeline and the only
component that operates entirely without an LLM. Its purpose is to transform the
incoming JSON Schema --- which may contain complex features such as \$ref
references, optional fields, validation constraints, and metadata
annotations --- into a normalised, decoder-safe form that minimises the risk of
state explosions in constrained decoding engines such as Outlines
\citep{willard2023efficient} and XGrammar~\citep{dong2025xgrammar}. The
importance of such schema preprocessing has been highlighted by recent work on
schema optimisation for LLM consumption, notably the PARSE framework
\citep{shrimal2025parse}.

\subsection{Operation Sequence}

The optimiser applies five transformations in a fixed pipeline:

\textbf{1. Inline resolution of \$ref references.} The input schema may contain
\$ref pointers in two forms: references to the \$defs section and references to
arbitrary paths within the schema. The optimiser resolves all such references
by replacing the \$ref key with the actual referenced sub-schema, traversing
the schema tree recursively. This produces a fully self-contained schema with
no external dependencies.

\textbf{2. Stripping of unsafe keys.} JSON Schema includes annotation and
validation keywords that are either irrelevant for generation or actively
harmful for constrained decoding. The optimiser removes the following keys from
every level of the schema: \texttt{default}, \texttt{title}, \texttt{format},
\texttt{minimum}, \texttt{maximum}, \texttt{minLength}, \texttt{maxLength},
\texttt{pattern}, \texttt{exclusiveMinimum}, \texttt{exclusiveMaximum},
\texttt{multipleOf}, and the entire \$defs section. The \texttt{format} key, in
particular, can cause decoders to allocate prohibitively large state machines
(for example, \texttt{format: date} may expand into a grammar covering
thousands of valid date strings). Stripping these keys reduces the average
schema size by approximately 30--40\% and eliminates a common class of decoder
failures.

\textbf{3. Enforcing \texttt{additionalProperties: false}.} The optimiser
inserts \texttt{additionalProperties: false} into every object node in the
schema. This is critical for constrained decoding: without it, the decoder's
finite-state machine must account for the possibility of arbitrary additional
keys at every object level, dramatically increasing the state space.

\textbf{4. Forcing all properties into \texttt{required}.} All properties at
every object level are moved into the \texttt{required} array. Fields that are
genuinely optional in the original schema are handled by wrapping their type
specification in an \texttt{anyOf} with \texttt{type: "null"}, ensuring that
the decoder permits \texttt{null} as a legal value while guaranteeing that
every key appears in the output. This prevents the decoder from omitting keys
and eliminates the non-determinism that arises when a decoder must decide
whether to emit a key or skip it at each optional field boundary.

\textbf{5. Extraction of FieldMetadata.} Finally, the optimiser traverses the
optimised schema to produce a list of \texttt{FieldMetadata} objects, each
containing: the field's dot-separated path, its JSON Schema type, its
\texttt{description} if present, its enum values if constrained, a flag
indicating optionality, and the nesting depth. This metadata is used downstream
by the DraftEngine to construct the \texttt{campo: valor} prompt and by the
adaptive pipeline's SchemaAnalyzer for task profiling.

\subsection{Impact on Downstream Decoders}

The importance of schema optimisation becomes apparent when constrained decoding
engines encounter unprocessed schemas. In preliminary experiments, we observed
that running a schema through the optimiser reduced grammar compilation time by
a factor of 3--5$\times$ and eliminated approximately 15\% of decoder-side
failures that would otherwise require fallback to the local parser. By investing
this optimisation upfront, the system ensures that the expensive LLM calls
operate on a schema guaranteed to be compatible with the downstream decoder,
shifting complexity from runtime error handling to a deterministic
preprocessing step.

% ==============================================================================
% 4. PRINCIPAL PIPELINE: DRAFT-THEN-DECODE
% ==============================================================================

\section{Principal Pipeline: Draft-then-Decode}
\label{sec:draft-decode}

The core contribution of the SEsml system is a \textbf{two-phase extraction
pipeline} that decouples semantic reasoning from structural formatting. This
section details the motivation behind the approach, the two processing phases
(DraftEngine and Constrained Decoding), and the post-validation layer that
ensures output reliability.

\subsection{Motivation}

Small Language Models under 14 billion parameters exhibit a phenomenon we term
\textbf{cognitive collapse} when tasked with simultaneous reasoning over Spanish
text and strict JSON syntax compliance. The model's probability distribution
must continuously arbitrate between linguistic localisation (resolving gendered
nouns, verb conjugations, and flexible word order in Spanish) and the rigid
token-level constraints of JSON (brackets, quotes, commas). This multitasking
overload manifests as omitted fields, malformed structures, and hallucinated
values.

The Draft-then-Decode architecture addresses this by splitting extraction from
formatting across two distinct LLM calls. This pattern is inspired by the
Generate \& Organise paradigm~\citep{li2024generate} and is closely related to
the Draft-Conditioned Constrained Decoding (DCCD) framework proposed
by~\citet{reddy2026dccd}, which provides a formal analysis of how draft
conditioning increases the feasible probability mass available to constrained
decoding, reducing the cumulative projection tax induced by hard structural
constraints. Our approach can be seen as an instantiation of this principle for
the specific domain of schema-guided information extraction from Spanish text,
where the semantic load of extraction and the structural load of JSON formatting
are particularly imbalanced.

\subsection{Phase 1: DraftEngine}

The DraftEngine is a single LLM call with \textbf{no structural constraints}.
The model receives the source text and target schema, and produces a
semi-structured intermediate representation composed of two XML-style tags:

\begin{itemize}
  \item \texttt{<Thinking>...</Thinking>}: A reasoning trace where the model
  performs field-by-field analysis, tracks evidence locations, and resolves
  ambiguities.
  \item \texttt{<Draft>campo: valor</Draft>}: One field per line in plain
  \texttt{campo: valor} (field: value) format. Values are plain text --- no
  JSON syntax, no quotes, no commas.
\end{itemize}

The prompt template used is version \textbf{v3b}, which evolved through three
major iterations during development. At approximately 1,163 characters (~387
tokens), it is designed to fit comfortably within the context windows of SLMs
while conveying all necessary constraints. The prompt is organised into five
sections demarcated by \texttt{===} delimiters: \texttt{TAREA} (task
description), \texttt{TEXTO} (input text), \texttt{ESQUEMA OBJETIVO} (optimised
schema), \texttt{INSTRUCCIONES} (extraction rules), and \texttt{FORMATO DE
RESPUESTA} (output template).

The \texttt{INSTRUCCIONES} section contains three critical rules:

\textbf{Rule 1: \texttt{[INFERENCIA]} --- Inference Control.} The model is
permitted to perform inference only when the schema field's \texttt{description}
contains specific trigger keywords: \texttt{reasoned}, \texttt{inferred},
\texttt{based on}, or \texttt{classify}. Inference is explicitly prohibited for
fields related to language, nationality, publication date, and external
metadata.

\textbf{Rule 2: \texttt{[VALORES]} --- Value Handling Protocol.} String values
must be copied exactly from the source text. Enum values must match the schema
definition case-sensitively. The value \texttt{null} is the only valid
representation of absence. For array fields, the model must enumerate all
instances exhaustively and return \texttt{[]} when zero instances are found.

\textbf{Rule 3: \texttt{[FORMATO]} --- Format Flexibility.} Rather than
imposing fixed formatting rules, each field's format is determined by its
\texttt{description} in the JSON Schema. This design choice keeps the prompt
schema-agnostic and defers format specification to the schema optimisation
phase.

A critical robustness mechanism: if the model's response lacks \texttt{<Draft>}
tags, the raw response text is passed directly to the decoder phase. The system
never discards a response.

\subsection{Phase 2: Constrained Decoding (Decoder)}

The second phase takes the draft produced by the DraftEngine and maps its
contents into a valid JSON object conforming to the optimised schema. Unlike the
DraftEngine, the Decoder call uses a strict
\texttt{response\_format=json\_schema} constraint, which forces the LLM to
generate only tokens that conform to the target JSON Schema.

The prompt for the Decoder is intentionally minimal --- approximately 150
characters: \emph{``Copy values from draft to output JSON. Output ONLY valid
JSON. Start directly with \{.''}

The core function \texttt{decode\_to\_json()} implements the following fallback
hierarchy:

\textbf{Attempt 1: Structured Output via
\texttt{response\_format=json\_schema}.} The primary path. In our experiments
with OpenAI-compatible endpoints serving SLMs, this succeeds in approximately
85--92\% of cases, depending on the model and schema complexity.

\textbf{Attempt 2: Local Fallback --- Direct Draft Parsing.} If the structured
output attempt fails, the system falls back to local parsing of the draft. The
\texttt{\_extract\_json()} helper applies four strategies in sequence:
(1) direct \texttt{json.loads}, (2) extraction from \texttt{```json} markdown
blocks, (3) bracket-matching to find the first \texttt{\{...\}} object, and
(4) parsing \texttt{campo: valor} lines into a flat dictionary.

\subsection{POST-Validation}

After either path produces a candidate JSON object, the \texttt{Verifier}
applies three deterministic validation layers:

\textbf{Layer 1: Structural Validation} checks the output against the schema for
type correctness, enum validity, and required field presence.

\textbf{Layer 2: Grounding Validation} checks every string value for presence in
the source text, using normalised comparison and a token-overlap fallback.

\textbf{Layer 3: Business Rules Validation} applies Pydantic model-level
constraints for inter-field dependencies.

\subsection{Why Draft-then-Decode Works}

The effectiveness of this two-phase architecture stems from a fundamental
property of autoregressive language models: the token-level probability
distribution cannot easily optimise for two conflicting objectives
simultaneously. As~\citet{reddy2026dccd} show, when an SLM must generate
outputs under hard structural constraints, standard constrained decoding
incurs a cumulative projection tax because the model assigns low probability
mass to valid continuations. The draft phase collapses the structural regime
(the model operates in a purely semantic mode), and the decode phase collapses
the semantic regime (the model performs a purely structural mapping). The draft
acts as a bridge representation: structured enough that the mapping to JSON is
deterministic in expectation, yet free enough that the model never needs to
simultaneously resolve Spanish syntax and JSON syntax. The local fallback path
provides a crucial safety net that guarantees output even when the constrained
decoding API fails entirely.

% ==============================================================================
% 5. ALTERNATIVE PIPELINES
% ==============================================================================

\section{Alternative Pipelines}
\label{sec:alternatives}

While the principal extraction pipeline achieves the highest accuracy through
its two-phase architecture, its reliance on two sequential LLM calls introduces
latency and cost that may be prohibitive in certain deployment scenarios. We
developed two alternative pipelines that compress the extraction process into a
single LLM call, each making different trade-offs between accuracy, latency,
and architectural complexity.

\subsection{HybridCoT: Single-Call with Visible Chain-of-Thought}

The \textbf{HybridCoTAgent} collapses the two-phase process into a single LLM
call by combining visible Chain-of-Thought reasoning~\citep{wei2022chain} with
direct JSON generation. The SchemaOptimizer normalises the input schema, after
which a single LLM call is issued with a prompt that instructs the model to
first reason within a \texttt{<Thinking>} block and then output the final JSON
directly. The JSON is extracted from the response using the
\texttt{\_extract\_json} four-strategy cascade.

The prompt template for this pipeline is nearly identical to the draft prompt
v3b but with two critical differences. First, the output format specification
requests \texttt{<Thinking>...</Thinking>} followed by a raw JSON object rather
than a \texttt{<Draft>} block. Second, the per-type value instructions are
expanded with explicit specifications for STRING, ENUM, NUMERIC, BOOLEAN, and
ARRAY fields, compensating for the absence of the
\texttt{response\_format=json\_schema} constraint.

HybridCoT achieves lower latency (one LLM call instead of two) and lower token
consumption, with a measured reduction of approximately 40--50\% in end-to-end
response time. However, it is less precise on complex schemas (complexity levels
L5--L10), where the absence of grammar-constrained decoding leads to higher
rates of malformed JSON and omitted fields. We recommend HybridCoT for simpler
schemas (L1--L4) and for latency-critical or cost-sensitive deployments.

\subsection{Adaptive Pipeline}

The \textbf{adaptive pipeline} is our most architecturally sophisticated
variant. It retains the single-call efficiency of HybridCoT but introduces a
retrieval-augmented prompt construction mechanism. The pipeline consists of
five stages:

\begin{enumerate}
  \item \textbf{SchemaOptimizer}: normalises the input schema.
  \item \textbf{SchemaAnalyzer}: builds a \texttt{TaskProfile} from the schema
  and task metadata, computing a 384-dimensional embedding via fastembed
  (\texttt{BAAI/bge-small-en-v1.5}) and retrieving the $K=15$ nearest
  neighbours from a precomputed development-set index.
  \item \textbf{PromptAssembler}: selects and orders modular prompt components
  based on the \texttt{TaskProfile} within a 1,400-character budget:
  \texttt{BASE} (always), \texttt{COMPLEXITY} (confidence $\geq 0.4$),
  \texttt{FIELD\_TYPES} (selective), \texttt{PATTERN\_ERRORS} (similarity
  $\geq 0.5$), \texttt{DOMAIN\_TIPS} (confidence $\geq 0.4$), and
  \texttt{FEW\_SHOT} (similarity $\geq 0.75$, not yet implemented).
  \item \textbf{Single LLM call}: the model receives the assembled prompt and
  produces visible CoT reasoning followed by direct JSON output.
  \item \textbf{JSON extraction} via the \texttt{\_extract\_json} cascade.
\end{enumerate}

The adaptive pipeline's key advantage is that domain-specific and
complexity-specific guidance demonstrably improves extraction accuracy on tasks
whose characteristics are well-covered by the development set. However, it
depends on precomputed embeddings and error maps; for genuinely novel tasks
whose representation falls far from any development exemplar, the assembled
prompt offers no advantage over HybridCoT.

% ==============================================================================
% 6. PROMPT ENGINEERING AND EVOLUTION
% ==============================================================================

\section{Prompt Engineering \& Evolution}
\label{sec:prompts}

The prompt templates used across all three SEsml pipelines underwent a sustained
engineering effort spanning four major versions. This section documents the
evolution from v1 to v3b (the active version), tracing the design reasoning
behind each iteration. The evolution reveals a counter-intuitive principle: for
SLMs, \textbf{shorter prompts consistently outperform longer, more detailed
ones} at the structured extraction task.

\subsection{Version History}

\textbf{v1 (deprecated).} At approximately 3,180 characters (~1,060 tokens),
the first prompt version was organised into four labelled blocks: \texttt{[A]}
for step-by-step instruction decomposition, \texttt{[B]} for a taxonomy of
evidence types, \texttt{[C]} for per-type formatting rules, and \texttt{[D]}
for a hierarchical conflict-resolution framework. The underlying intuition was
that more detail would produce more reliable extraction. In practice, at 1,060
tokens the prompt consumed more than a quarter of the context window of a 4K
SLM, leading to instruction truncation and what we term \emph{instruction
dilution}: the model's attention, distributed across a lengthy prompt, assigned
lower effective weight to critical formatting constraints.

\textbf{v2 (deprecated).} In response, v2 was cut to approximately 1,551
characters (~517 tokens), a reduction of over 50\%. The four-block structure
was abandoned in favour of a single continuous text. However, at approximately
170 words of instruction beyond the task description, v2 was still too long for
the smallest SLMs, which omitted \texttt{<Draft>} tags entirely.

\textbf{v3 (deprecated).} The third version restructured the prompt around
``3 critical rules'' --- the precursor to INFERENCIA, VALORES, and FORMATO ---
but ironically made the prompt longer by introducing a structured preamble and
more explicit examples. In several configurations, the combined
prompt-plus-input exceeded 4,000 tokens, causing context overflow.

\textbf{v3b (active).} At approximately 1,163 characters (~387 tokens), v3b
represents the sweet spot balancing completeness and brevity. Three key design
decisions distinguish it:

\begin{enumerate}
  \item \textbf{Elimination of the four-block structure.} The three rules are
  presented as flat imperative statements, saving ~200 characters without
  reducing instructional content.
  \item \textbf{Three clear imperative rules using Spanish keywords.} Each rule
  uses a Spanish keyword header (\texttt{[INFERENCIA]}, \texttt{[VALORES]},
  \texttt{[FORMATO]}) as a categorical anchor. Maintaining Spanish throughout
  the prompt improves compliance on Spanish-specific linguistic patterns.
  \item \textbf{Schema-driven formatting.} Instead of prescribing format rules
  in the prompt, v3b defers all formatting decisions to the schema's
  \texttt{description} field. This keeps the prompt schema-agnostic and saves
  approximately 300 characters.
\end{enumerate}

\subsection{Design Principles}

The iterative refinement crystallised into four general design principles:

\textbf{1. Brevity beats completeness.} SLMs --- particularly quantized 4-bit
and 8-bit variants --- follow short, imperative instructions more reliably than
detailed, explanatory ones. Reducing the prompt from ~1,060 to ~387 tokens did
not reduce instruction adherence; it improved it.

\textbf{2. Rules, not explanations.} The most effective formulations state
\textbf{what} the model should do without explaining \textbf{why}. Explanations
consume tokens without adding operational constraints.

\textbf{3. Task-language alignment.} Maintaining Spanish throughout the
prompt --- even for structural instructions --- improved compliance on
Spanish-specific linguistic patterns.

\textbf{4. Schema-driven formatting.} Format instructions should be part of the
schema, not the prompt. This separation of concerns means that formatting
improvements benefit all three pipelines simultaneously.

\subsection{Constrained Decoding Prompt}

The decoder prompt underwent a similar but more dramatic evolution from ~400
characters to the current ~150-character version:
\emph{``Copy values from draft to output JSON. Output ONLY valid JSON. Start
directly with \{.''} This reduction was possible because
\texttt{response\_format=json\_schema} already guarantees syntactic JSON
validity, making explicit format instructions redundant.

% ==============================================================================
% 7. RESULTS
% ==============================================================================

\section{Results}
\label{sec:results}

At the time of writing, the official GenSIE 2026 evaluation results have not
been released by the organisers. The analysis in this section is framed in terms
of the expected behaviour of each pipeline given its design characteristics,
along with a template for reporting the official results once available.

\subsection{Evaluation Setup}

The GenSIE 2026 shared task evaluates participant systems in a zero-shot,
model-agnostic setting. Each submitted pipeline is evaluated against a suite of
evaluation models comprising both published open-weight language models and a
set of held-out models. The models are constrained to the SLM regime
(approximately 14 billion parameters and below) and include representative
architectures such as Llama 3.1 8B, Qwen2.5 7B, Gemma 3, and other comparable
models.

The primary evaluation metric is \textbf{GapClosed}, defined as the fraction of
the error gap between a baseline system and perfect performance that a
participant system eliminates. Formally, if $E_{\text{baseline}}$ is the error
rate of a reference baseline and $E_{\text{system}}$ is the error rate of the
submitted system, then:

\[
\text{GapClosed} = \frac{E_{\text{baseline}} - E_{\text{system}}}{E_{\text{baseline}}}
\]

where error is measured using a flattened Micro-F1 score that treats each leaf
field in the JSON Schema as an independent binary classification decision. A
GapClosed of 1.0 indicates perfect performance, while 0.0 indicates performance
equivalent to the baseline.

Secondary metrics include \textbf{Raw Micro-F1} and \textbf{Efficiency}, defined
as Micro-F1 divided by total token consumption.

\subsection{Expected Outcomes}

Although numerical results are pending, the architectural properties of each
pipeline allow us to formulate expected performance patterns.

\textbf{Extraction pipeline (principal).} We expect the two-phase draft-then-
decode pipeline to achieve the highest overall GapClosed, particularly on
schemas of complexity levels L5--L10, where the constrained decoding step
guarantees structural validity and the strict inference rules reduce
hallucination.

\textbf{Hybrid\_cot pipeline.} We anticipate this pipeline will be most
competitive on simpler schemas (L1--L4), achieving the highest Efficiency score
among the three pipelines. On complex schemas, the absence of constrained
decoding may lead to increased rates of malformed output.

\textbf{Adaptive pipeline.} The adaptive pipeline's performance is expected to
be more variable and domain-dependent. On tasks whose instruction embeddings
have high similarity to development-set profiles, it may approach or match the
extraction pipeline's accuracy while consuming only a single LLM call. On novel
schemas, it offers no advantage over HybridCoT.

% ==============================================================================
% 8. ERROR ANALYSIS
% ==============================================================================

\section{Error Analysis}
\label{sec:errors}

\subsection{DraftEngine Failures}

The DraftEngine, operating without structural constraints, is the most volatile
component. We identified four recurrent failure patterns:

\textbf{Missing structural tags.} In approximately 5--10\% of cases,
particularly with models at 3B--7B parameters, the response lacks either
\texttt{<Thinking>} or \texttt{<Draft>} tags. The fail-open design mitigates
this by passing the raw response to the decoder.

\textbf{Malformed \texttt{campo: valor} lines.} Common defects include missing
colons, multi-line values, and unescaped colons. Recovery is imperfect but
partially handled by the fallback parser.

\textbf{Truncation.} The model may stop generating before all fields are listed,
especially for schemas with 15+ fields or 4+ nesting levels. The downstream
decoder maps only the available subset.

\textbf{Repetition loops.} In approximately 1--3\% of cases, the model enters a
reasoning loop within the \texttt{<Thinking>} block, consuming output tokens and
crowding out the \texttt{<Draft>} section.

\subsection{Decoder Failures}

\textbf{Structured output API errors.} Some SLMs raise errors when the schema
contains large enum sets (50+ values) or deeply nested \texttt{anyOf} chains.
The SchemaOptimiser eliminates many triggers, but edge cases remain.

\textbf{Draft-independent generation.} In approximately 3--5\% of cases, the
decoder ignores the draft-as-context and generates a JSON object based solely
on the schema definition. The resulting output is structurally valid but
semantically incorrect.

\textbf{Empty or null outputs.} Models below 3B parameters occasionally return
empty responses, handled by the local fallback.

\subsection{Common SLM Error Classes}

Across all three pipelines, six recurrent error classes appear:

\begin{enumerate}
  \item \textbf{Enum mismatch.} SLMs produce values that do not case-match the
  schema definition (e.g., \texttt{"Positive"} vs. \texttt{"POSITIVE"}).
  \item \textbf{Date format violations.} Models default to natural-language
  date expressions instead of ISO format.
  \item \textbf{Null hallucination.} Models fill missing fields with plausible
  default values instead of \texttt{null}.
  \item \textbf{Entity span boundary errors.} Models truncate or extend named
  entity spans.
  \item \textbf{Array omissions.} Models enumerate only a subset of instances
  for array fields.
  \item \textbf{Language interference.} Models occasionally emit values in
  English despite Spanish prompts and input.
\end{enumerate}

\subsection{Adaptive-Pipeline-Specific Errors}

The adaptive pipeline introduces additional failure modes related to its
retrieval-augmented prompt assembly: low-similarity profiles (task falls far
from development data), domain cue irrelevance (incorrect or too broad domain
labels), and error map generalisation failure (held-out schemas exhibit error
modes not present in development data).

% ==============================================================================
% 9. CONCLUSIONS
% ==============================================================================

\section{Conclusions}
\label{sec:conclusions}

We have presented SEsml, a system for zero-shot schema-guided information
extraction from Spanish text using Small Language Models. The system's core
architectural insight is the decoupling of semantic reasoning from structural
formatting through a two-phase draft-then-decode pipeline, which mitigates the
cognitive collapse that SLMs experience when asked to simultaneously reason
about linguistic content and enforce JSON syntax compliance. This approach is
theoretically grounded in the analysis of~\citet{reddy2026dccd} and
practically instantiated through the Generate \& Organise paradigm
\citep{li2024generate}.

The draft-then-decode principle forms the basis of our principal submission,
complemented by two alternative pipelines offering different trade-offs along
the accuracy--latency--adaptability axes. The HybridCoT pipeline collapses the
two phases into a single LLM call with visible Chain-of-Thought
reasoning~\citep{wei2022chain}, while the adaptive pipeline introduces
retrieval-augmented prompt construction through the SchemaAnalyzer and
PromptAssembler components.

Our development experience yielded several practical lessons for prompt
engineering with SLMs. First, brevity and imperative rules consistently
outperform detailed explanatory instructions for models under 14 billion
parameters. Second, schema optimisation is not merely a preprocessing
convenience but a critical enabler for constrained decoding engines. Third, the
separation of concerns between a reasoning draft and a formatting decoder proved
robust to a wide range of schema complexities and model architectures,
suggesting that the draft-then-decode pattern may generalise beyond information
extraction to other structured generation problems with SLMs.

\textbf{Future work.} Several extensions are planned. The full few-shot
integration in the adaptive pipeline --- currently disabled with the similarity
threshold defined but not triggering example inclusion --- would allow retrieval
of relevant exemplar pairs. A multi-model routing layer could select the
appropriate pipeline dynamically based on task complexity estimates. Automatic
post-correction using grounding validation feedback offers a path to recover
from common error classes without full retry loops. Finally, cross-lingual
evaluation beyond Spanish would test the generalisation of the draft-then-decode
principle across Romance languages.

The SEsml system is available as open-source
software\footnote{\url{https://github.com/alejbv/SEsml.git}}, and all three
pipelines can be reproduced using the provided Docker containers and evaluation
harness.

% ==============================================================================
% ACKNOWLEDGEMENTS
% ==============================================================================

\section*{Acknowledgements}

We thank the organisers of the GenSIE 2026 shared task at IberLEF 2026 for
providing the evaluation framework and dataset. Computational resources were
provided by the Faculty of Mathematics and Computing, Universidad de La Habana.

% ==============================================================================
% REFERENCES
% ==============================================================================

\bibliography{references}

\end{document}
